\section{The Limits of Today’s AI}
\label{sec:limits}

Modern AI can recognise faces, translate languages, and even drive cars.
But under the hood, these systems are \emph{statistical mirrors}: they
reflect the patterns that humans have already put into the data.

\begin{tcolorbox}[colback=blue!5, colframe=blue!40, title=Quick thought]
If you train an AI on 10\,000 labelled photos of cats and dogs, it
becomes very good at telling cats from dogs.  Yet it has no idea what
a cat \emph{is} – it only knows that a certain pixel pattern usually
comes with the label “cat”.
\end{tcolorbox}

Three problems arise from this approach:
\begin{enumerate}
    \item \textbf{Every concept must be predefined.}  The AI cannot invent
    a new category that a human did not already think of.
    \item \textbf{Time is treated as a fixed step.}  The AI does not
    experience the world as a flow of events.
    \item \textbf{Bias is baked in.}  If the training data contains
    stereotypes, the AI will reproduce them.
\end{enumerate}

\begin{tcolorbox}[colback=red!5, colframe=red!40, title=The Parrot and the Scientist]
Think of today’s large language models as a \textbf{parrot}.  The parrot
hears the word “water” thousands of times and learns to say it in the
right context.  But it has never touched water, never tasted it, never
felt its weight.  It has no idea what water really \emph{is}.

Now imagine a \textbf{scientist} who has never heard the word “water”.
She studies a puddle: she measures its temperature, watches how it
evaporates, sees how it reflects light.  She builds a concept —
“that transparent, flowing thing that disappears in sunlight” — and
gives it her own name.  Her understanding is grounded in observation,
not in borrowed vocabulary.

Apeiron is the scientist.  It does not repeat what humans have said;
it builds knowledge from the ground up.
\end{tcolorbox}

\begin{tcolorbox}[colback=red!5, colframe=red!40, title=Why Apeiron is fundamentally different]
Apeiron does \emph{not} try to make a better version of today’s AI.
It is a completely different path — one that starts with raw sensory
data and lets structure emerge on its own, without any human‑provided
labels, categories, or clocks.  Think of it not as a faster car, but
as the first spaceship.
\end{tcolorbox}